\chapter{行业痛点与设计目标}

\section{痛点一：日志供给已经超过集中式分析的承载边界}

日志是工业设备、操作系统、中间件和生产服务最贴近故障现场的数字证据，但其增长速度已使“先集中、后分析”面临现实瓶颈。Bloomberg 在 USENIX LISA 2019 公布的生产实践中，其遥测基础设施从数万台机器与应用采集数据，每秒处理超过 600 万项指标、每天产生约 120 TB 日志\cite{terng2019bloomberg}。更早的系统研究给出了集中搬运的量级推演：按每台机器仅产生 1 MB/s 日志计算，1000 节点集群一天即可产生约 86 TB 数据；若希望在 48 分钟内完成集中传输，需要约 300 Gb/s 的专用网络与磁盘带宽\cite{logothetis2011insitu}。这说明日志分析的首要矛盾并非“有没有模型”，而是原始数据增长速度、现场响应时限与可用传输资源之间的不匹配。

日志不仅规模大，而且高度异构。Loghub-2.0 收录的 14 类大规模标注数据集平均每类包含约 360 万行日志和 249 个模板，其中 Spark 与 Thunderbird 数据分别超过 1600 万行，Thunderbird 包含 1241 个模板\cite{jiang2024loghub2}。不同设备、软件版本和开发团队使用不同格式与语义，同一故障又可能以错误码、超时、重试或性能退化等不同症状出现。因而，单纯增加集中存储容量只能解决“收得到”，不能自动解决“看得懂、来得及、处置得了”。

更关键的是，日志知识并非静态不变。EvLog 对真实软件演化数据的研究发现，版本演进会同时带来日志解析错误、日志事件演化和事件序列不稳定三类问题\cite{huo2023evlog}。这意味着依赖固定关键词、固定模板或一次性训练分布的方案，会随着版本升级和现场工况变化发生能力衰减；恰恰是未被既有知识覆盖的新模板和异常组合，往往更需要被保留并交给更强的语义分析能力。

上述事实形成了典型的价值长尾：高频、重复、确定性较强的事件适合由低成本方法快速收敛；数量较少但新颖、高风险、低置信或跨节点关联的事件，则需要更大的上下文窗口和推理能力。若所有日志采用相同处理深度，系统不是在简单事件上浪费计算与通信资源，就是在复杂事件上牺牲诊断质量。

\section{痛点二：分布式故障的症状在局部，根因在关联关系中}

工业现场与云原生基础设施都具有多节点、多层依赖和异步交互特征。一个存储节点的抖动，可能先表现为上游服务请求超时，继而触发连接池耗尽、任务重试和队列堆积；单条日志记录的通常是本节点观察到的症状，而非传播链起点。USENIX 对真实 HBase/Hadoop 生态的研究将诊断难点归纳为分布式日志不一致、有效信息被噪声淹没以及根因判断具有概率性，并指出日志分散和缺少生态级上下文会直接增加排障难度\cite{li2013hadoopdiagnosis}。另一项 NSDI 研究通过关联日志事件，在三个真实分布式系统中定位了六类性能问题，说明根因识别依赖跨组件事件关系，而不仅是孤立关键词\cite{nagaraj2012distalyzer}。

这带来两项工程约束。第一，系统必须在采集阶段统一时间、来源、设备、模板、严重度和事件标识，否则跨节点关联缺少共同坐标。第二，全局分析不能等同于把全部历史日志无差别塞入大模型：冗余上下文既增加 Token、时延和费用，也会稀释真正相关的证据。更合理的路径是先在边缘形成事件级初判，再围绕少量高价值主事件按时间邻近、来源关系、模板相似性和严重度筛选上下文，把有限的云端推理预算用于回答“共同根因是什么、故障如何传播、下一步如何处置”。

\section{痛点三：云端能力更强，但不能成为现场业务的单点依赖}

云端模型具有更大的参数规模和更强的开放语义能力，却同时受到网络时延、服务限流、调用预算、数据边界和 Provider 可用性的约束。Uptime Institute 2026 年中断分析显示，57\% 的受访者表示最近一次重大中断成本超过 10 万美元，20\% 超过 100 万美元；与光纤和连接相关的外部基础设施中断正在增加，而且更容易造成长时间影响\cite{uptime2026outage}。对要求连续运行的工业现场而言，这意味着实时告警主链不能以“网络可用且云端按时返回”为前提。

常见的固定级联方案通常在模型置信度低于某个阈值时上云，但“是否值得上云”与“此刻能否上云”是两个不同问题：前者取决于事件风险、边缘认知不确定性、新颖度和可能影响范围，后者取决于网络健康、队列压力、调用预算、限流和隐私策略。只使用单一置信度阈值，会把业务价值判断与系统可用性判断混为一谈；简单失败重试或先进先出缓存，又可能使链路恢复后低价值请求挤占高风险事件的诊断机会。

因此，云边协同不应只是把两个模型串联起来，而应解决一个更本质的问题：\textbf{在连续事件流和动态资源约束下，何时、为何、以多大代价调用哪一级智能能力。}这是一项事件级认知资源调度问题，也是 CoTriage 区别于“边缘分类器加云端接口”的核心出发点。

\begin{table}[htbp]
  \centering
  \footnotesize
  \renewcommand{\arraystretch}{0.96}
  \caption{现场日志诊断的核心矛盾与 CoTriage 的设计响应}
  \label{tab:problem-response}
  \begin{tabularx}{\textwidth}{>{\bfseries}p{0.19\textwidth} X X}
    \toprule
    核心矛盾 & 传统路径的边界 & CoTriage 的设计响应 \\
    \midrule
    数据规模与处理时限 & 全量集中传输占用带宽，深度分析难以覆盖连续高频事件 & 就近事件化与边缘早停，仅对高价值事件按需上云 \\
    模板知识与持续演化 & 固定规则和静态模型难覆盖新模板、版本漂移与组合异常 & 规则、模板、轻量分类器和边缘小模型组成可治理分层推理链 \\
    局部症状与全局根因 & 单节点日志难还原共同依赖、传播顺序和集群级影响 & 围绕主事件筛选跨来源关联上下文，云端执行结构化深诊 \\
    事件价值与资源约束 & 固定阈值无法同时表达风险、网络、队列、预算与策略 & RAISE 分离“值得发送”与“当前可发送”，输出三动作路由 \\
    业务连续与弱网波动 & 云端链路失败会阻塞判断，普通缓存又缺少风险优先级 & 本地判断持续输出，Outbox 持久化并按风险优先重放 \\
    \bottomrule
  \end{tabularx}
\end{table}

\section{设计目标：构建事件级、可约束、可演进的云边协同系统}

CoTriage 不替代现有日志平台，也不单独追求某个分类模型的离线得分，而是把多节点采集、边缘判断、资源决策、云端诊断、反馈写回和配置治理组织成可运行闭环。统一事件使每条日志都能回答：边缘掌握了什么证据，为什么选择本地、云端或延期路径，结果如何被追踪并反哺决策。系统形成六项设计要求：

{\small
\begin{enumerate}[leftmargin=*,itemsep=1pt,topsep=3pt,parsep=0pt]
  \item \textbf{现场实时自治。} 确定性事件在边缘快速收敛；云端不可达时仍持续产生基础判断，不让网络故障转化为感知链路停摆。
  \item \textbf{分析能力按价值分配。} 以风险、不确定性、新颖度、历史边云分歧和跨设备影响估计云端增益，将深度推理集中用于困难或高风险事件。
  \item \textbf{资源状态进入决策闭环。} 将网络、队列、预算、限流、隐私和运维策略纳入动态门槛，使路由随事件与环境共同变化。
  \item \textbf{全局诊断有证据边界。} 围绕主事件聚合有限且相关的跨节点上下文，在 Token 预算内输出根因候选、处置建议和结构化解释。
  \item \textbf{延期任务可靠且有优先级。} 对暂不可发送但仍具价值的事件进行持久化、幂等重试和死信治理，恢复后优先处理高风险事件。
  \item \textbf{全链可观测、可治理、可扩展。} 记录事件、推理路径、版本、路由依据和诊断状态，以稳定接口支持来源、模型、Provider 与节点扩展。
\end{enumerate}
}
